\chapter{Overview}

\section{Scientific context}
Iron homeostasis is tightly regulated at the cellular level. The models studied
in this project describe key species (labile iron, ferritin and iron storage,
IRP1, TfR1, ferroportin, and later HIF) and their kinetic interactions under
different iron-replenishment conditions. The numerical goal is to sample a
high-dimensional parameter space, retain only trajectories that satisfy
biological constraints, and characterise the resulting admissible parameter
regions.

\section{Software lineage}
The original simulation framework was developed by Eric Fanchon and Alexandre
Rocca (TIMC-IMAG, CNRS). Irina Piazza later reorganised the code for clarity and
maintainability: configuration was centralised, input/output was standardised
with pandas DataFrames, filtering and logging were rewritten, plotting utilities
were added, and the ODE right-hand side was unified. The public release
distributed with this documentation is the HIF--iron oriented version of that
work, modernised for Python~3.

Changes introduced by Irina Piazza are marked in the source with comments of the
form \texttt{\# IP <date>}.

\section{Architecture of the public code}
The main modules are:

\begin{description}
  \item[\texttt{simulate.py}] Entry point. Implements sampling modes
  (\texttt{Box-Sampling}, \texttt{Re-Sampling}, \texttt{Shell-Sampling},
  \texttt{PCA-guided-sampling}, \ldots), numerical integration, constraint
  checks, optional filtering, dumping of trajectories, and plotting.
  \item[\texttt{model.py}] Model-specific ODE system, algebraic steady-state
  construction, absolute/relative constraints, and optional stiffness
  diagnostics.
  \item[\texttt{config.py}] Global constants: solver tolerances, sampling
  settings, variable/parameter names, bounds and biological constraints.
  \item[\texttt{filter.py}] Post-processing filter on ratios and concentration
  bounds; writes filtered solution lists and enclosing boxes.
  \item[\texttt{read\_file.py}] Reads bound files and results tables.
  \item[\texttt{log.py}] Writes \texttt{info.txt} and copies input files into
  each run directory.
  \item[\texttt{sobol\_lib.py}, \texttt{pca.py}] Quasi-Monte Carlo sampling and
  PCA-guided resampling helpers.
  \item[\texttt{plot\_*.py}, \texttt{multiplot\_*.py}] Offline visualisation
  utilities.
\end{description}

Each execution creates a timestamped folder \texttt{Run\_DD\_MM\_YYYY-HhMmSs/}
containing results, plots, filtered outputs (when applicable), and run metadata.

\section{Typical pipeline}
\begin{enumerate}
  \item Sample parameter vectors inside a box (Sobol sequence) or around
  previous solutions (re-sampling / shell / PCA modes).
  \item For each vector, compute algebraic initial conditions, improve the
  steady state numerically if needed, then integrate the ODEs for each iron
  input listed in \texttt{MULTIPLE\_INPUTS}.
  \item Accept the set only if absolute and relative constraints are satisfied.
  \item After the sampling loop, filter the accepted set with
  \texttt{filter.py}, dump a subset of trajectories, and produce plots.
\end{enumerate}
